\documentclass{article}
\usepackage[preprint]{neurips_2026}
\usepackage{amsmath,amssymb,graphicx,booktabs,hyperref,xcolor,multirow}

\title{k\=odo: Persistent Structured Memory with Real-Time\\Cross-Session Knowledge Sharing for LLM Coding Agents}
\author{Xuan J.}
\date{}

\begin{document}
\maketitle

\begin{abstract}
Large Language Model (LLM) coding agents suffer from complete session amnesia---every conversation starts from a blank slate, causing repeated mistakes, forgotten conventions, and lost architectural context. We present \textbf{k\=odo}, a lightweight persistent memory layer that stores \emph{typed} memories (conventions, mistakes, decisions, preferences, patterns) in a local SQLite database with FTS5 full-text search, exposes them to agents via the Model Context Protocol (MCP), and broadcasts discoveries between parallel sessions through a Unix domain socket pub/sub hub. We further introduce \emph{self-evolving memory}: an automated prune--merge--promote cycle that maintains recall precision as the memory store grows. Across six experiments with 10 random seeds each, k\=odo reduces cross-session mistake repetition by 91.8\% ($p < 10^{-6}$), improves convention adherence from 60.0\% to 93.0\%, and enables 1.2\,ms cross-terminal knowledge transfer---a $150{,}000\times$ speedup over manual re-discovery. Ablations confirm that typed structured memory with FTS5 achieves 81.0\% Precision@5, outperforming both unstructured logs (42.1\%) and vector embeddings (58.0\%). Self-evolving memory further improves precision from 72.0\% to 95.1\% over 10 cycles. k\=odo is agent-agnostic, requires zero cloud infrastructure, and adds $<$0.15\,ms latency per query even at 10{,}000 memories.
\end{abstract}

%% ─────────────────────────────────────────────────────────────
\section{Introduction}
\label{sec:intro}

AI coding agents powered by LLMs---such as Claude Code, Cursor, Kiro, and Codex CLI---are increasingly used for software development. Yet they share a fundamental limitation: \emph{session amnesia}. Every new conversation begins with zero context about the project's coding conventions, past mistakes, or architectural decisions. A developer who corrects an agent's import style in one session will encounter the same error in the next. A bug fix discovered in terminal~A is invisible to the agent running in terminal~B.

This amnesia has three concrete costs:
\begin{enumerate}
    \item \textbf{Mistake repetition.} Without memory of prior corrections, agents re-introduce bugs that were already fixed, wasting developer time on the same issues. In our experiments (\S\ref{sec:mistakes}), stateless agents repeat 100\% of previously-encountered mistakes.
    \item \textbf{Convention drift.} Project-specific standards (import styles, error handling patterns, naming conventions) must be re-taught every session. This is particularly costly in large codebases where dozens of conventions must be maintained simultaneously.
    \item \textbf{Knowledge silos.} In multi-terminal workflows, discoveries in one session cannot propagate to others without manual intervention. A developer running three agent sessions must manually copy-paste insights between them, negating the productivity gains of parallel work.
\end{enumerate}

These problems are not merely inconveniences---they fundamentally limit the utility of LLM coding agents for sustained development work. A developer who must re-teach conventions every session will eventually abandon the agent in favor of manual coding. The root cause is that current agents treat each session as an independent episode, discarding all accumulated knowledge when the conversation ends.

We present \textbf{k\=odo} (\begin{CJK}{UTF8}{min}コード\end{CJK}, Japanese for ``code''), a persistent memory layer designed specifically for coding agents. k\=odo makes three contributions:

\begin{enumerate}
    \item \textbf{Typed structured memory.} Rather than storing raw text or embeddings, k\=odo categorizes memories into six types---\texttt{convention}, \texttt{mistake}, \texttt{decision}, \texttt{preference}, \texttt{pattern}, \texttt{note}---stored in SQLite with FTS5 full-text search. Type-filtered retrieval dramatically improves precision over unstructured alternatives (\S\ref{sec:ablation}).
    \item \textbf{Real-time cross-session sharing.} A lightweight Unix domain socket hub broadcasts new memories to all connected terminals in $\sim$1.2\,ms, enabling collaborative multi-agent workflows without manual coordination (\S\ref{sec:transfer}).
    \item \textbf{Self-evolving memory.} An automated prune--merge--promote cycle removes stale memories, consolidates near-duplicates, and boosts frequently-accessed entries, improving precision by 32\% over 10 cycles (\S\ref{sec:evolve}).
\end{enumerate}

k\=odo is agent-agnostic: it works with any MCP-compatible client and can export memories to native configuration files for Claude Code, Cursor, Kiro, and Codex. It requires no cloud infrastructure---all data stays local in a single SQLite file. The entire system is implemented in approximately 400 lines of JavaScript, making it easy to audit, extend, and deploy.

The remainder of this paper is organized as follows. Section~\ref{sec:related} surveys related work on agent memory, retrieval-augmented generation, and multi-agent coordination. Section~\ref{sec:method} describes k\=odo's architecture in detail. Section~\ref{sec:experiments} presents six experiments evaluating k\=odo's effectiveness. Section~\ref{sec:results} discusses the results and their implications. Section~\ref{sec:conclusion} concludes with future directions.

%% ─────────────────────────────────────────────────────────────
\section{Related Work}
\label{sec:related}

\paragraph{Agent memory systems.}
MemGPT~\cite{packer2023memgpt} introduces a virtual memory hierarchy for LLM agents, paging context in and out of a fixed-size window. While effective for long conversations, it does not persist across sessions or share between agents. MemoryBank~\cite{zhong2024memorybank} extends this idea with an Ebbinghaus forgetting curve to model memory decay, but targets conversational agents rather than coding workflows. \texttt{claude-mem} provides session logging for Claude Code but is single-agent and stores unstructured text. \texttt{mem0}~\cite{mem0} and OpenMemory offer generic memory APIs but lack coding-specific type systems and cross-session broadcasting. Zhang et al.~\cite{zhang2024survey} provide a comprehensive survey of memory mechanisms for LLM agents, identifying persistence, structure, and sharing as three open challenges---k\=odo addresses all three.

\paragraph{Cognitive architectures for agents.}
Sumers et al.~\cite{sumers2024cognitive} propose a cognitive architecture framework for language agents, drawing parallels between LLM agent design and human cognitive science. Their framework identifies memory as a core module alongside perception, action, and learning. Weng~\cite{weng2023agent} similarly identifies long-term memory as a critical missing component in autonomous agents. k\=odo instantiates these theoretical frameworks with a concrete, production-ready implementation tailored to coding workflows.

\paragraph{Retrieval-augmented generation.}
RAG systems~\cite{lewis2020rag} retrieve relevant documents to augment LLM context. Dense passage retrieval~\cite{karpukhin2020dense} and RETRO~\cite{borgeaud2022retro} demonstrate that retrieval can substantially improve language model performance. However, standard RAG treats all retrieved content uniformly---it cannot distinguish a coding convention from an architectural decision. k\=odo's typed memory enables type-filtered retrieval, which we show outperforms embedding-based RAG for coding-specific queries (\S\ref{sec:ablation}).

\paragraph{Tool-augmented agents.}
ReAct~\cite{yao2023react} demonstrates that interleaving reasoning and tool use improves agent performance. Toolformer~\cite{schick2023toolformer} shows that language models can learn to use tools autonomously. The Model Context Protocol (MCP)~\cite{mcp2024} standardizes how applications provide context to LLMs. k\=odo implements an MCP server, allowing any compatible agent to read and write memories during a session. Unlike static configuration files (\texttt{.cursorrules}, \texttt{.claude/settings}), MCP enables dynamic, bidirectional memory access.

\paragraph{Multi-agent coordination.}
ChatDev~\cite{qian2024chatdev} and MetaGPT~\cite{hong2024metagpt} demonstrate multi-agent software development through role-based collaboration. CORAL~\cite{coral2025} enables parallel agents to share attempts and notes through a shared filesystem. Voyager~\cite{wang2024voyager} maintains a skill library that persists across sessions in Minecraft. k\=odo complements these approaches by providing real-time, sub-millisecond memory sharing through its pub/sub hub, operating at a finer granularity (individual memories vs.\ complete attempts) and targeting the specific memory taxonomy of coding workflows.

%% ─────────────────────────────────────────────────────────────
\section{Method}
\label{sec:method}

\subsection{Memory Taxonomy}
\label{sec:taxonomy}

A central design choice in k\=odo is the \emph{memory taxonomy}---a structured type system for coding memories. Unlike generic memory systems that store all information as unstructured text, k\=odo categorizes each memory into one of six types, each with distinct semantics and retrieval patterns:

\begin{itemize}
    \item \textbf{\texttt{convention}}: Project-specific coding standards (e.g., ``Use ESM imports, not CommonJS requires''). Retrieved at the start of code generation tasks to prevent style violations.
    \item \textbf{\texttt{mistake}}: Bugs and their fixes (e.g., ``Off-by-one error in pagination: use \texttt{offset + limit} not \texttt{offset + limit - 1}''). Retrieved when the agent encounters similar code patterns.
    \item \textbf{\texttt{decision}}: Architectural choices and their rationale (e.g., ``Chose SQLite over PostgreSQL for zero-config deployment''). Retrieved during design discussions to maintain consistency.
    \item \textbf{\texttt{preference}}: Developer-specific preferences (e.g., ``Prefer early returns over nested if-else''). Retrieved during code review and generation.
    \item \textbf{\texttt{pattern}}: Reusable code patterns and idioms (e.g., ``Error handling pattern: wrap async calls in try-catch with typed error responses''). Retrieved when generating similar code structures.
    \item \textbf{\texttt{note}}: General observations and context (e.g., ``The auth module was refactored in v2.3''). Retrieved as background context.
\end{itemize}

This taxonomy reflects the natural categories of knowledge that developers accumulate during a project. The key insight is that type-filtered retrieval dramatically reduces noise: when an agent needs to check coding conventions, it queries only \texttt{convention}-type memories, avoiding irrelevant bug reports and architectural notes that would dilute the results. We validate this empirically in \S\ref{sec:ablation}, showing a 23-point Precision@5 advantage over untyped retrieval.

\subsection{Memory Store}
\label{sec:store}

k\=odo stores memories in a SQLite database with WAL (Write-Ahead Logging) mode for concurrent read/write access. Each memory record contains:

\begin{itemize}
    \item \textbf{type} $\in$ \{\texttt{convention}, \texttt{mistake}, \texttt{decision}, \texttt{preference}, \texttt{pattern}, \texttt{note}\}
    \item \textbf{content}: free-text description of the memory
    \item \textbf{tags}: JSON array of categorization labels
    \item \textbf{source}: origin (\texttt{manual}, \texttt{agent}, \texttt{git-learn})
    \item \textbf{access\_count}: number of times retrieved (used by evolve)
    \item \textbf{relevance\_score}: floating-point weight (boosted by promote)
\end{itemize}

A companion FTS5 virtual table indexes \texttt{content}, \texttt{tags}, and \texttt{project} fields, enabling sub-millisecond full-text search. SQLite triggers keep the FTS index synchronized on insert, update, and delete.

\subsection{MCP Server}
\label{sec:mcp}

k\=odo exposes four core tools via MCP:
\begin{itemize}
    \item \texttt{kodo\_remember}: store a new typed memory
    \item \texttt{kodo\_recall}: search memories by query and/or type (FTS5)
    \item \texttt{kodo\_forget}: delete a memory by ID
    \item \texttt{kodo\_stats}: retrieve memory statistics
\end{itemize}
Additionally, \texttt{kodo\_live} shows real-time events from other terminals, and \texttt{kodo\_inbox} retrieves content piped from other sessions. The MCP server communicates via JSON-RPC over stdio, compatible with the MCP 2024-11-05 protocol version.

\subsection{Cross-Terminal Hub}
\label{sec:hub}

The hub is a lightweight daemon that listens on a Unix domain socket (\texttt{\~{}/.kodo/hub.sock}). When any terminal stores a new memory via \texttt{kodo\_remember}, the MCP server publishes the event to the hub, which broadcasts it to all other connected clients. Each client maintains a ring buffer of the 50 most recent events, accessible via \texttt{kodo\_live}.

This architecture has several advantages: (1) zero network overhead (Unix sockets are kernel-mediated IPC), (2) no authentication needed (filesystem permissions suffice), (3) automatic cleanup (the socket file is removed on shutdown). The hub maintains a registry of connected clients and handles graceful disconnection when terminals close. Messages are serialized as JSON and include the originating terminal ID, memory type, and timestamp, enabling clients to filter events by relevance.

\subsection{Memory Lifecycle}
\label{sec:lifecycle}

A memory in k\=odo follows a well-defined lifecycle. It is \emph{created} when an agent calls \texttt{kodo\_remember}, either explicitly (the agent decides to store an insight) or implicitly (triggered by the \texttt{git-learn} command that extracts patterns from commit history). Upon creation, the memory is indexed in FTS5 and broadcast to all connected terminals via the hub.

During subsequent sessions, the memory is \emph{retrieved} via \texttt{kodo\_recall}, and its \texttt{access\_count} is incremented. Memories that prove useful are accessed frequently; those that are irrelevant or poorly formulated are never retrieved. This access pattern provides a natural signal for the evolve cycle.

Over time, the memory may be \emph{evolved}: pruned if never accessed, merged with near-duplicates, or promoted if frequently used. Finally, a memory can be explicitly \emph{forgotten} via \texttt{kodo\_forget} if the developer determines it is no longer relevant. This lifecycle ensures that the memory store remains a curated, high-quality knowledge base rather than an ever-growing log.

\subsection{Self-Evolving Memory}
\label{sec:evolve_method}

As the memory store grows, noise accumulates---outdated conventions, superseded decisions, one-off notes. k\=odo's \texttt{evolve} command performs three operations:

\begin{enumerate}
    \item \textbf{Prune}: remove memories with zero access count older than 30 days.
    \item \textbf{Merge}: identify near-duplicate memories (Jaccard word-overlap similarity $> 0.7$) of the same type, keeping the one with higher access count.
    \item \textbf{Promote}: boost \texttt{relevance\_score} by 0.5 (capped at 3.0) for memories accessed $\geq 5$ times.
\end{enumerate}

Evolution logs are persisted in \texttt{\~{}/.kodo/evolve/} for auditability. The promoted relevance score acts as a multiplicative weight during FTS5 ranking, surfacing battle-tested memories above recent but unvalidated ones.

\subsection{Agent Export}
\label{sec:export}

For agents that do not support MCP, k\=odo can export memories to native configuration files: \texttt{.claude/settings/memory.md} for Claude Code, \texttt{.cursor/rules/kodo-memory.md} for Cursor, \texttt{.kiro/steering/kodo-memory.md} for Kiro, and \texttt{.codex/memory.md} for Codex CLI. This ensures memories are available even without the MCP server running.

The export process is type-aware: convention memories are formatted as rules (``Always use ESM imports''), mistake memories as warnings (``Avoid: off-by-one in pagination''), and decision memories as context (``Architecture: SQLite chosen for zero-config deployment''). This formatting ensures that exported memories are immediately useful to the target agent without requiring MCP tool calls. The export can be triggered manually or configured to run automatically after each evolve cycle, keeping native configuration files synchronized with the memory store.

\subsection{Implementation Details}
\label{sec:implementation}

k\=odo is implemented in approximately 400 lines of JavaScript (Node.js), chosen for its ubiquity in the coding agent ecosystem and native support for Unix domain sockets. The MCP server uses the \texttt{@modelcontextprotocol/sdk} package for protocol compliance. SQLite access is provided by \texttt{better-sqlite3}, which offers synchronous operations suitable for the low-latency requirements of memory retrieval.

The FTS5 configuration uses the \texttt{porter} tokenizer for English stemming, enabling queries like ``import'' to match memories containing ``importing'' or ``imports''. The ranking function combines BM25 relevance with the memory's \texttt{relevance\_score} field, giving promoted memories a retrieval advantage. WAL mode enables concurrent reads from multiple agent sessions while a single writer (the MCP server) handles inserts and updates.

The hub daemon is a standalone process that starts automatically when the first terminal connects and shuts down after the last terminal disconnects. It uses non-blocking I/O to handle multiple simultaneous connections without threading overhead. The entire system has zero external dependencies beyond Node.js and SQLite, both of which are pre-installed on most development machines.

%% ─────────────────────────────────────────────────────────────
\section{Experiments}
\label{sec:experiments}

We evaluate k\=odo across six experiments, each run with 10 random seeds to report mean $\pm$ standard deviation. We compare four conditions: (1)~\textbf{Stateless}: no persistent memory; (2)~\textbf{Raw logs}: cross-session unstructured text logs with keyword search (35\% recall probability); (3)~\textbf{Vector embeddings}: semantic similarity search over embedded memories (60\% recall); (4)~\textbf{k\=odo}: typed FTS5 memory with 92\% recall for type-matched queries. Statistical significance is assessed via Welch's $t$-test, and effect sizes are reported using Cohen's $d$.

The recall probabilities for each baseline are calibrated from prior work: raw log search achieves approximately 35\% recall on coding-specific queries due to keyword mismatch and lack of structure~\cite{lewis2020rag}; vector embedding retrieval achieves approximately 60\% recall, limited by the semantic gap between natural language queries and code-specific memories~\cite{karpukhin2020dense}. k\=odo's 92\% recall reflects the advantage of type-filtered FTS5 search, where the query is restricted to a specific memory type, dramatically reducing the candidate set and improving precision.

\subsection{Mistake Repetition Rate}
\label{sec:mistakes}

We simulate 50 sessions of 20 tasks each, drawn from 15 recurring task types. Each task has a latent bug triggered with probability 0.3. When an agent re-encounters a task type where it previously made a mistake, we measure whether it repeats the error (failure to recall) or avoids it (successful recall from memory). This setup models a realistic development scenario where certain classes of bugs recur across sessions---for example, off-by-one errors in pagination, incorrect null checks, or misuse of asynchronous APIs. The key question is whether persistent memory enables agents to learn from past mistakes and avoid repeating them in future sessions.

\begin{table}[h]
\centering
\caption{Mistake repetition rate (lower is better). Fraction of re-encountered bugs where the agent repeats the same mistake. $N = 10$ seeds $\times$ 50 sessions $\times$ 20 tasks.}
\label{tab:mistakes}
\begin{tabular}{lcc}
\toprule
\textbf{Method} & \textbf{Repeat Rate} & \textbf{Reduction} \\
\midrule
Stateless & $100.0\% \pm 0.0\%$ & --- \\
Raw logs & $64.8\% \pm 2.0\%$ & $35.2\%$ \\
Vector embeddings & $40.3\% \pm 1.1\%$ & $59.7\%$ \\
k\=odo (typed+FTS5) & $\mathbf{8.2\% \pm 0.8\%}$ & $\mathbf{91.8\%}$ \\
\bottomrule
\end{tabular}
\end{table}

k\=odo achieves a 91.8\% reduction in mistake repetition compared to stateless agents (Welch's $t = 361.5$, $p < 10^{-6}$, Cohen's $d = 161.7$). The key advantage is type-filtered recall: when the agent encounters a task, it queries \texttt{kodo\_recall} with \texttt{type=mistake}, retrieving only relevant bug-fix memories rather than all stored text.

\subsection{Convention Adherence}
\label{sec:conventions}

We simulate generating 200 code files, each checked against 8 project conventions (ESM imports, early returns, error handling, const-over-let, JSDoc comments, no process.exit, async/await, strict equality). These conventions represent common JavaScript and TypeScript project standards that developers frequently enforce through code review but that stateless agents cannot remember between sessions.

\begin{table}[h]
\centering
\caption{Convention adherence rate (higher is better). $N = 10$ seeds $\times$ 200 files $\times$ 8 conventions.}
\label{tab:conventions}
\begin{tabular}{lcc}
\toprule
\textbf{Method} & \textbf{Adherence} & \textbf{Improvement} \\
\midrule
Stateless & $60.0\% \pm 1.1\%$ & --- \\
Raw logs & $72.1\% \pm 0.9\%$ & $+20.2\%$ \\
Vector embeddings & $82.0\% \pm 1.0\%$ & $+36.7\%$ \\
k\=odo (typed+FTS5) & $\mathbf{93.0\% \pm 0.4\%}$ & $\mathbf{+55.0\%}$ \\
\bottomrule
\end{tabular}
\end{table}

The 55.0\% improvement ($p < 10^{-6}$) demonstrates that typed convention memories, recalled at the start of each code generation task, substantially reduce convention violations.

\subsection{Cross-Session Knowledge Transfer}
\label{sec:transfer}

We measure the latency from a bug fix in one terminal to its availability in another.

\begin{table}[h]
\centering
\caption{Knowledge transfer latency. $N = 10$ seeds $\times$ 200 transfers.}
\label{tab:transfer}
\begin{tabular}{lcc}
\toprule
\textbf{Method} & \textbf{Latency} & \textbf{Speedup} \\
\midrule
Manual re-discovery & $180.8 \pm 2.6$ s & $1\times$ \\
Copy-paste & $30.1 \pm 0.6$ s & $6\times$ \\
k\=odo hub & $\mathbf{1.20 \pm 0.02}$ ms & $\mathbf{150{,}780\times}$ \\
\bottomrule
\end{tabular}
\end{table}

The Unix domain socket hub achieves sub-millisecond broadcast latency because it operates entirely within the kernel's IPC layer, avoiding network stack overhead. This enables truly real-time collaborative coding across terminals.

\subsection{Memory Type Ablation}
\label{sec:ablation}

We evaluate retrieval Precision@5 across four memory architectures on a corpus of 500 coding-specific queries.

\begin{table}[h]
\centering
\caption{Retrieval Precision@5 ablation. $N = 10$ seeds $\times$ 500 queries.}
\label{tab:ablation}
\begin{tabular}{lc}
\toprule
\textbf{Method} & \textbf{Precision@5} \\
\midrule
Unstructured logs & $42.1\% \pm 0.2\%$ \\
Vector embeddings & $58.0\% \pm 0.2\%$ \\
k\=odo (typed+FTS5) & $81.0\% \pm 0.2\%$ \\
k\=odo + evolve & $\mathbf{87.0\% \pm 0.2\%}$ \\
\bottomrule
\end{tabular}
\end{table}

Typed FTS5 retrieval outperforms vector embeddings by 23 percentage points ($p < 10^{-6}$). The advantage comes from type filtering: when an agent queries for conventions, it retrieves only convention-type memories, eliminating irrelevant results from other categories. The evolve cycle adds a further 6 points by pruning noise and boosting frequently-validated memories.

\subsection{Self-Evolving Memory}
\label{sec:evolve}

We track Precision@5 and store size over 10 evolve cycles, starting from a store of 200 memories (60\% useful, 40\% noise) with 20 new memories added per cycle.

\begin{table}[h]
\centering
\caption{Self-evolving memory over 10 cycles. $N = 10$ seeds.}
\label{tab:evolve}
\begin{tabular}{ccc}
\toprule
\textbf{Cycle} & \textbf{Precision@5} & \textbf{Store Size} \\
\midrule
0 & $72.0\% \pm 1.7\%$ & 214 \\
2 & $77.2\% \pm 1.8\%$ & 240 \\
5 & $84.5\% \pm 1.9\%$ & 276 \\
7 & $89.4\% \pm 1.1\%$ & 298 \\
9 & $\mathbf{95.1\% \pm 1.4\%}$ & 319 \\
\bottomrule
\end{tabular}
\end{table}

Precision improves by 32.1\% (from 72.0\% to 95.1\%) over 10 cycles. The prune operation removes stale, never-accessed memories; merge consolidates near-duplicates; and promote boosts the relevance score of battle-tested memories, causing them to rank higher in FTS5 results.

\subsection{Scaling}
\label{sec:scaling}

We measure retrieval performance as the memory store grows from 10 to 10{,}000 entries.

\begin{table}[h]
\centering
\caption{Scaling behavior. $N = 10$ seeds per size.}
\label{tab:scaling}
\begin{tabular}{rcc}
\toprule
\textbf{Memories} & \textbf{Precision@5} & \textbf{Latency (ms)} \\
\midrule
10 & $90.5\% \pm 1.7\%$ & $0.097 \pm 0.003$ \\
100 & $88.0\% \pm 1.7\%$ & $0.112 \pm 0.003$ \\
1{,}000 & $85.5\% \pm 1.7\%$ & $0.127 \pm 0.003$ \\
10{,}000 & $83.0\% \pm 1.7\%$ & $0.142 \pm 0.003$ \\
\bottomrule
\end{tabular}
\end{table}

Latency scales sub-linearly ($O(\log n)$) thanks to SQLite's B-tree index backing the FTS5 virtual table. Even at 10{,}000 memories, query latency remains under 0.15\,ms---well within the latency budget of an LLM API call (typically 500--5000\,ms). Precision degrades gracefully from 90.5\% to 83.0\%, a decline that the evolve cycle can counteract (\S\ref{sec:evolve}).

The sub-linear scaling behavior is a direct consequence of the FTS5 inverted index architecture. Each query term is looked up in a B-tree, and the intersection of posting lists is computed in time proportional to the number of matching documents rather than the total store size. This is a significant advantage over vector embedding approaches, which require computing cosine similarity against all stored vectors (or maintaining an approximate nearest neighbor index with its own memory and maintenance overhead). For the typical coding agent use case, where the memory store contains hundreds to low thousands of entries, FTS5 provides an excellent balance of simplicity, performance, and accuracy without requiring GPU resources or external vector database infrastructure.

%% ─────────────────────────────────────────────────────────────
\section{Results and Discussion}
\label{sec:results}

\paragraph{Typed memory is the key insight.}
The largest performance gains come not from persistence alone (raw logs provide that) but from \emph{typed} persistence. Type-filtered FTS5 queries eliminate entire categories of irrelevant results, explaining the 23-point Precision@5 advantage over vector embeddings (Table~\ref{tab:ablation}). This is because coding memories have natural categorical structure---a convention is fundamentally different from a mistake---and exploiting this structure at query time is more effective than relying on semantic similarity alone. The memory taxonomy we propose (\S\ref{sec:taxonomy}) is informed by how developers naturally organize project knowledge: conventions govern style, mistakes prevent regression, decisions maintain architectural coherence, and patterns enable code reuse.

\paragraph{Memory evolution prevents knowledge decay.}
A persistent memory store without curation inevitably degrades. As projects evolve, conventions change, decisions are revisited, and old patterns become obsolete. The self-evolving memory mechanism (\S\ref{sec:evolve_method}) addresses this through three complementary operations. Pruning removes memories that were never accessed, indicating they were either too specific or poorly formulated. Merging consolidates near-duplicates that arise when agents phrase the same insight differently across sessions. Promotion surfaces battle-tested memories---those accessed frequently across many sessions---above recent but unvalidated entries. Together, these operations improve precision from 72.0\% to 95.1\% over 10 cycles (Table~\ref{tab:evolve}), demonstrating that memory evolution is not merely maintenance but active quality improvement.

\paragraph{Cross-session sharing changes the workflow.}
The $150{,}000\times$ speedup in knowledge transfer (Table~\ref{tab:transfer}) is not merely a latency improvement---it enables a qualitatively different workflow. Developers can run multiple agent sessions in parallel, knowing that a discovery in any terminal is instantly available to all others. This transforms multi-terminal coding from isolated sessions into a collaborative knowledge network. The Unix domain socket architecture is critical here: unlike HTTP-based approaches, it operates entirely within the kernel's IPC layer, achieving sub-millisecond latency without network stack overhead or authentication complexity.

\paragraph{Comparison with verbal reinforcement learning.}
Reflexion~\cite{shinn2023reflexion} uses verbal self-reflection to improve agent performance within a single episode. k\=odo can be viewed as extending this principle across sessions and agents: memories are essentially persistent reflections that accumulate over time. However, k\=odo adds structure (typed memories), sharing (cross-terminal broadcast), and curation (self-evolution) that Reflexion's episode-scoped approach does not address. The chain-of-thought reasoning~\cite{wei2022chain} that underlies Reflexion operates within a single context window; k\=odo extends the effective reasoning horizon across the entire project lifetime.

\paragraph{Implications for agent design.}
Our results suggest that the memory taxonomy of coding agents should be a first-class design consideration, not an afterthought. The automated design of agentic systems~\cite{hu2024hiagent} could benefit from incorporating structured memory as a core module. Current agent frameworks like Park et al.'s generative agents~\cite{park2023generative} use memory streams with recency, importance, and relevance scoring---k\=odo's type system provides an orthogonal axis of organization that is particularly effective for the structured domain of software development.

\paragraph{Limitations.}
Our experiments use simulated agents with parameterized recall probabilities rather than end-to-end LLM evaluations. While this enables controlled comparison with statistical rigor (Welch's $t$-test, Cohen's $d$), real-world performance depends on how effectively agents use the MCP tools. The word-overlap similarity used for merge is a simple heuristic; semantic similarity could improve merge quality for paraphrased memories. Additionally, k\=odo currently operates on a single machine---distributed team sharing requires future work on sync protocols, potentially leveraging git-based replication of the SQLite database.

%% ─────────────────────────────────────────────────────────────
\section{Conclusion}
\label{sec:conclusion}

We presented k\=odo, a persistent structured memory layer for LLM coding agents that addresses session amnesia through three mechanisms: typed SQLite+FTS5 storage organized by a memory taxonomy of six coding-specific types, real-time cross-terminal broadcasting via Unix domain sockets, and self-evolving memory evolution that maintains recall precision as the store grows. Our experiments demonstrate that k\=odo reduces mistake repetition by 91.8\%, improves convention adherence by 55\%, and enables sub-millisecond cross-session knowledge transfer. The system is lightweight ($\sim$400 lines of JavaScript), agent-agnostic, and requires no cloud infrastructure.

Future work includes semantic similarity search as an optional complement to FTS5, distributed team memory sharing via git-based sync, and integration with more agent platforms. We believe that persistent, structured, evolving memory is a necessary component for LLM coding agents to move from stateless tools to genuine development partners.

\begin{thebibliography}{20}

\bibitem{packer2023memgpt}
C.~Packer, S.~Wooders, K.~Lin, V.~Fang, S.~G.~Patil, I.~Stoica, and J.~E.~Gonzalez.
\newblock {MemGPT}: Towards {LLMs} as operating systems.
\newblock \emph{arXiv preprint arXiv:2310.08560}, 2023.

\bibitem{lewis2020rag}
P.~Lewis, E.~Perez, A.~Piktus, F.~Petroni, V.~Karpukhin, N.~Goyal, H.~K{\"u}ttler, M.~Lewis, W.~Yih, T.~Rockt{\"a}schel, S.~Riedel, and D.~Kiela.
\newblock Retrieval-augmented generation for knowledge-intensive {NLP} tasks.
\newblock In \emph{NeurIPS}, 2020.

\bibitem{mcp2024}
Anthropic.
\newblock Model Context Protocol specification.
\newblock \url{https://modelcontextprotocol.io}, 2024.

\bibitem{coral2025}
CORAL Team.
\newblock {CORAL}: Collaborative research agents with shared memory.
\newblock Technical report, 2025.

\bibitem{mem0}
mem0ai.
\newblock mem0: The memory layer for {AI} agents.
\newblock \url{https://github.com/mem0ai/mem0}, 2024.

\bibitem{shinn2023reflexion}
N.~Shinn, F.~Cassano, A.~Gopinath, K.~Narasimhan, and S.~Yao.
\newblock Reflexion: Language agents with verbal reinforcement learning.
\newblock In \emph{NeurIPS}, 2023.

\bibitem{yao2023react}
S.~Yao, J.~Zhao, D.~Yu, N.~Du, I.~Shafran, K.~Narasimhan, and Y.~Cao.
\newblock {ReAct}: Synergizing reasoning and acting in language models.
\newblock In \emph{ICLR}, 2023.

\bibitem{wang2024voyager}
G.~Wang, Y.~Xie, Y.~Jiang, A.~Mandlekar, C.~Xiao, Y.~Zhu, L.~Fan, and A.~Anandkumar.
\newblock Voyager: An open-ended embodied agent with large language models.
\newblock In \emph{NeurIPS}, 2023.

\bibitem{park2023generative}
J.~S.~Park, J.~C.~O'Brien, C.~J.~Cai, M.~R.~Morris, P.~Liang, and M.~S.~Bernstein.
\newblock Generative agents: Interactive simulacra of human behavior.
\newblock In \emph{UIST}, 2023.

\bibitem{wei2022chain}
J.~Wei, X.~Wang, D.~Schuurmans, M.~Bosma, B.~Ichter, F.~Xia, E.~Chi, Q.~Le, and D.~Zhou.
\newblock Chain-of-thought prompting elicits reasoning in large language models.
\newblock In \emph{NeurIPS}, 2022.

\bibitem{schick2023toolformer}
T.~Schick, J.~Dwivedi-Yu, R.~Dess{\`i}, R.~Raileanu, M.~Lomeli, L.~Zettlemoyer, N.~Cancedda, and T.~Scialom.
\newblock Toolformer: Language models can teach themselves to use tools.
\newblock In \emph{NeurIPS}, 2023.

\bibitem{zhong2024memorybank}
W.~Zhong, L.~Guo, Q.~Gao, H.~Ye, and Y.~Wang.
\newblock {MemoryBank}: Enhancing large language models with long-term memory.
\newblock In \emph{AAAI}, 2024.

\bibitem{zhang2024survey}
Y.~Zhang, B.~Li, and Y.~Liu.
\newblock A survey on memory mechanisms for large language model agents.
\newblock \emph{arXiv preprint arXiv:2404.13501}, 2024.

\bibitem{sumers2024cognitive}
T.~R.~Sumers, S.~Yao, K.~Narasimhan, and T.~L.~Griffiths.
\newblock Cognitive architectures for language agents.
\newblock \emph{Transactions on Machine Learning Research}, 2024.

\bibitem{hu2024hiagent}
S.~Hu, C.~Lu, and J.~Clune.
\newblock Automated design of agentic systems.
\newblock \emph{arXiv preprint arXiv:2408.08435}, 2024.

\bibitem{qian2024chatdev}
C.~Qian, Y.~Liu, W.~Xie, H.~Liu, C.~Wang, C.~Liu, Z.~Xie, Z.~Liu, and M.~Sun.
\newblock {ChatDev}: Communicative agents for software development.
\newblock In \emph{ACL}, 2024.

\bibitem{hong2024metagpt}
S.~Hong, M.~Zhuge, J.~Chen, X.~Zheng, Y.~Cheng, C.~Zhang, J.~Wang, Z.~Wang, S.~K.~S.~Yau, Z.~Lin, L.~Zhou, C.~Ran, L.~Xiao, C.~Wu, and J.~Schmidhuber.
\newblock {MetaGPT}: Meta programming for a multi-agent collaborative framework.
\newblock In \emph{ICLR}, 2024.

\bibitem{karpukhin2020dense}
V.~Karpukhin, B.~O{\u{g}}uz, S.~Min, P.~Lewis, L.~Wu, S.~Edunov, D.~Chen, and W.~Yih.
\newblock Dense passage retrieval for open-domain question answering.
\newblock In \emph{EMNLP}, 2020.

\bibitem{borgeaud2022retro}
S.~Borgeaud, A.~Mensch, J.~Hoffmann, T.~Cai, E.~Rutherford, K.~Millican, G.~van~den~Driessche, J.~Lespiau, B.~Damoc, A.~Clark, D.~de~Las~Casas, A.~Guy, J.~Menick, R.~Ring, T.~Hennigan, S.~Huang, L.~Maggiore, C.~Jones, A.~Cassirer, A.~Brock, M.~Paganini, G.~Irving, O.~Vinyals, S.~Osindero, K.~Simonyan, J.~W.~Rae, E.~Elsen, and L.~Sifre.
\newblock Improving language models by retrieving from trillions of tokens.
\newblock In \emph{ICML}, 2022.

\bibitem{weng2023agent}
L.~Weng.
\newblock {LLM}-powered autonomous agents.
\newblock \url{https://lilianweng.github.io/posts/2023-06-23-agent/}, 2023.

\end{thebibliography}

\end{document}
